\chapter{Posterior and Explanation Summaries}
\label{app:explainability}\label{app:explanation}

\section{Prior-Specific Posterior Quantities}
Table~\ref{tab:gen-prior-signals} describes the scale of the continuous variance-mixing weights and horseshoe shrinkage proxies in the first repeat and first fold. Their definitions differ, while the coefficient and hazard-ratio quantities share an interpretation across all three priors. Continuous effects are expressed per training-fold standard deviation, and categorical effects compare the encoded level with its reference category.

% Generated from hash-verified saved evidence. Do not edit by hand.
\begin{table}[htbp]
\centering
\caption{Illustrative prior-specific summaries from the first outer fit.}
\label{tab:gen-prior-signals}
\small
\setlength{\tabcolsep}{3pt}
\begin{adjustbox}{max width=\textwidth}
\begin{tabular}{llrrrr}
\toprule
Cohort & Prior & Encoded features & Minimum & Median & Maximum \\
\midrule
GBSG2 & Spike-and-slab & 8 & 0.355 & 0.379 & 0.417 \\
METABRIC & Horseshoe & 19 & 0.169 & 0.590 & 0.751 \\
METABRIC & Spike-and-slab & 19 & 0.342 & 0.397 & 0.484 \\
WHAS500 & Horseshoe & 9 & 0.096 & 0.665 & 0.677 \\
WHAS500 & Spike-and-slab & 9 & 0.354 & 0.368 & 0.501 \\
\bottomrule
\end{tabular}
\end{adjustbox}
\par\smallskip
\begin{minipage}{\textwidth}\footnotesize Continuous spike-and-slab entries describe variance-mixing weights. Horseshoe entries describe a local-scale shrinkage proxy. These prior-specific summaries have different definitions. Examples use the first repeat and first fold.\end{minipage}
\end{table}


\section{Comparison with CoxPH Explanations}
The common intrinsic importance measure is the posterior expected absolute contribution of a clinical predictor group to the linear predictor. Signed encoded contributions are summed within a group before taking the absolute value. Table~\ref{tab:gen-explanation-cross-model} compares these rankings with exact linear SHAP for the paired CoxPH model. Its differences use the same folds and explanation seeds.

% Generated from hash-verified saved evidence. Do not edit by hand.
\begin{table}[htbp]
\centering
\caption{Secondary comparison with explanations of the paired CoxPH model.}
\label{tab:gen-explanation-cross-model}
\small
\setlength{\tabcolsep}{3pt}
\begin{adjustbox}{max width=\textwidth}
\begin{tabular}{llrrr}
\toprule
Cohort & Bayesian prior & Folds & Cox agreement $\rho$ & Same-model minus Cox $\rho$ [95\% interval] \\
\midrule
GBSG2 & Horseshoe & 7/25 & 0.810 & +0.156 [+0.068, +0.256] \\
GBSG2 & Normal & 25/25 & 0.925 & +0.002 [-0.011, +0.015] \\
GBSG2 & Spike-and-slab & 25/25 & 0.922 & +0.007 [-0.010, +0.025] \\
METABRIC & Horseshoe & 23/25 & 0.909 & -0.042 [-0.060, -0.026] \\
METABRIC & Normal & 25/25 & 0.868 & -0.004 [-0.016, +0.006] \\
METABRIC & Spike-and-slab & 25/25 & 0.870 & -0.008 [-0.019, +0.002] \\
WHAS500 & Horseshoe & 25/25 & 0.859 & -0.001 [-0.043, +0.042] \\
WHAS500 & Normal & 25/25 & 0.871 & +0.009 [-0.016, +0.035] \\
WHAS500 & Spike-and-slab & 25/25 & 0.871 & +0.002 [-0.028, +0.028] \\
\bottomrule
\end{tabular}
\end{adjustbox}
\par\smallskip
\begin{minipage}{\textwidth}\footnotesize Cox agreement compares Bayesian intrinsic contributions with exact linear SHAP for the paired CoxPH risk score. Differences use matched folds and explanation seeds. Intervals are descriptive hierarchical-bootstrap estimates.\end{minipage}
\end{table}


\section{Leading Predictor Groups}
Table~\ref{tab:gen-explanation-top-groups} presents the three groups with the largest mean posterior expected absolute contributions for each cohort and prior. Mean rank, top-five frequency and resampling intervals describe how their positions vary. This display summarises fitted model behaviour; all predictor groups remain part of their specified models.

% Generated from hash-verified saved evidence. Do not edit by hand.
\clearpage
\begingroup
\singlespacing
\setlength{\parskip}{0pt}
\setlength{\parindent}{0pt}
\footnotesize
\renewcommand{\arraystretch}{1.03}
\setlength{\tabcolsep}{3pt}
\setlength{\LTpre}{6pt}
\setlength{\LTpost}{6pt}
\begin{longtable}{>{\raggedright\arraybackslash}p{2cm}>{\raggedright\arraybackslash}p{3cm}>{\raggedright\arraybackslash}p{4.6cm}>{\raggedright\arraybackslash}p{1.4cm}>{\raggedright\arraybackslash}p{3.5cm}}
\caption{Leading intrinsic contribution groups, with descriptive rank uncertainty.}\label{tab:gen-explanation-top-groups}\\
\toprule
Cohort & Prior & Predictor group & \shortstack[l]{Mean\\rank} & \shortstack[l]{Top-five frequency\\95\% interval} \\
\midrule
\endfirsthead
\multicolumn{5}{l}{\tablename\ \thetable\ (continued)}\\
\toprule
Cohort & Prior & Predictor group & \shortstack[l]{Mean\\rank} & \shortstack[l]{Top-five frequency\\95\% interval} \\
\midrule
\endhead
\midrule
\multicolumn{5}{r}{Continued on next page}\\
\endfoot
\bottomrule
\endlastfoot
GBSG2 & Normal & Progesterone receptor & 1.11 & 1.000 [1.000, 1.000] \\
GBSG2 & Normal & Positive lymph nodes & 1.96 & 1.000 [1.000, 1.000] \\
GBSG2 & Normal & Hormonal therapy & 4.16 & 0.880 [0.720, 1.000] \\
GBSG2 & Horseshoe & Progesterone receptor & 1.43 & 1.000 [1.000, 1.000] \\
GBSG2 & Horseshoe & Positive lymph nodes & 1.57 & 1.000 [1.000, 1.000] \\
GBSG2 & Horseshoe & Tumour grade & 3.29 & 1.000 [1.000, 1.000] \\
GBSG2 & Spike-and-slab & Progesterone receptor & 1.11 & 1.000 [1.000, 1.000] \\
GBSG2 & Spike-and-slab & Positive lymph nodes & 1.93 & 1.000 [1.000, 1.000] \\
GBSG2 & Spike-and-slab & Hormonal therapy & 4.16 & 0.880 [0.720, 1.000] \\
WHAS500 & Normal & Age & 1.00 & 1.000 [1.000, 1.000] \\
WHAS500 & Normal & Heart rate & 2.20 & 1.000 [1.000, 1.000] \\
WHAS500 & Normal & Diastolic blood pressure & 3.40 & 1.000 [1.000, 1.000] \\
WHAS500 & Horseshoe & Age & 1.00 & 1.000 [1.000, 1.000] \\
WHAS500 & Horseshoe & Heart rate & 2.12 & 1.000 [1.000, 1.000] \\
WHAS500 & Horseshoe & Diastolic blood pressure & 3.36 & 1.000 [1.000, 1.000] \\
WHAS500 & Spike-and-slab & Age & 1.00 & 1.000 [1.000, 1.000] \\
WHAS500 & Spike-and-slab & Heart rate & 2.20 & 1.000 [1.000, 1.000] \\
WHAS500 & Spike-and-slab & Diastolic blood pressure & 3.36 & 1.000 [1.000, 1.000] \\
METABRIC & Normal & Age at diagnosis & 1.00 & 1.000 [1.000, 1.000] \\
METABRIC & Normal & Nottingham prognostic index & 2.11 & 1.000 [1.000, 1.000] \\
METABRIC & Normal & ER status & 3.16 & 1.000 [1.000, 1.000] \\
METABRIC & Horseshoe & Age at diagnosis & 1.00 & 1.000 [1.000, 1.000] \\
METABRIC & Horseshoe & Nottingham prognostic index & 2.00 & 1.000 [1.000, 1.000] \\
METABRIC & Horseshoe & ER status & 3.67 & 1.000 [1.000, 1.000] \\
METABRIC & Spike-and-slab & Age at diagnosis & 1.00 & 1.000 [1.000, 1.000] \\
METABRIC & Spike-and-slab & Nottingham prognostic index & 2.08 & 1.000 [1.000, 1.000] \\
METABRIC & Spike-and-slab & ER status & 3.17 & 1.000 [1.000, 1.000] \\
\end{longtable}
\noindent\footnotesize Three groups per cohort and prior are displayed by mean posterior expected absolute contribution. Frequencies and intervals use hierarchical resampling. The global sampling-stability assessment is reported separately in Table~\ref{tab:gen-explanation-gate}.\par
\endgroup

